% ===========================================================================
% Section: Conclusion
% ===========================================================================
\section{Conclusion}
\label{sec:conclusion}

We have presented \winlic, a Windows-compatible offline multi-modal
Gaussian Splatting framework that eliminates the ROS dependency of prior
Gaussian SLAM systems. The framework makes three key contributions:

\begin{enumerate}
    \item A \textbf{pluggable sensor source interface} (\texttt{ISensorSource})
    that decouples the reconstruction pipeline from data acquisition,
    supporting both offline file-based datasets and live Orbbec camera
    streams through a uniform four-method interface. The design is
    extensible to future sensor SDKs without pipeline modification.

    \item A \textbf{multi-modal depth fusion} module that combines sparse
    LiDAR depth, dense iToF depth, and stereo SGM depth via
    confidence-weighted averaging, producing higher-quality depth maps
    than any single modality. The fused depth is further refined by the
    SPNet TensorRT completion engine.

    \item A \textbf{ROS-free, cross-platform build} that runs on both
    Windows (MSVC) and Linux (GCC) with pure CMake, preserving the core
    Gaussian Splatting algorithm of \gaussianlic{} without quality
    degradation.
\end{enumerate}

\subsection{Limitations and Future Work}

The current framework has several limitations that we plan to address:

\begin{itemize}
    \item \textbf{Pose dependency:} The system relies on pre-computed
    poses from an external SLAM front-end (e.g., FAST-LIVO2). Integrating
    a lightweight LIO frontend would enable fully self-contained
    processing, though at the cost of additional platform dependencies
    (PCL, IKFOM).

    \item \textbf{Orbbec SDK integration:} The current implementation
    provides a stub for the Orbbec SDK. Full integration with live
    streaming and on-device point-cloud generation requires testing
    with physical hardware.

    \item \textbf{GUI:} The current CLI interface is sufficient for
    batch processing but lacks interactive visualization. A Qt or
    Dear~ImGui-based GUI would improve usability for parameter tuning
    and real-time preview.

    \item \textbf{Large-scale scenes:} The in-memory Gaussian
    representation does not support out-of-core processing. For
    city-scale scenes, a submap-based tiling strategy (similar to
    Co-LIC2) would be needed.
\end{itemize}

\subsection{Reproducibility}

The complete source code, build instructions, and configuration files
are provided in the \texttt{win\_lic/} directory of the project repository.
The original ROS-based code in \texttt{src/} is preserved untouched,
allowing direct comparison between the two implementations.
